\documentclass[11pt,a4paper]{article}

\usepackage[T1]{fontenc}
\usepackage[utf8]{inputenc}
\usepackage{charter}
\usepackage[scaled=0.9]{helvet}
\usepackage{microtype}
\usepackage[margin=2.4cm]{geometry}
\usepackage{xcolor}
\usepackage{booktabs,tabularx}
\usepackage{enumitem}
\usepackage{titlesec}
\usepackage{fancyhdr}
\usepackage{listings}
\usepackage{needspace}
\usepackage[most]{tcolorbox}
\usepackage[colorlinks=true,linkcolor=accent,urlcolor=accent]{hyperref}

\definecolor{accent}{HTML}{1F4E79}
\definecolor{codebg}{HTML}{F4F6F8}
\definecolor{warnbg}{HTML}{FFF4E5}
\definecolor{warnframe}{HTML}{C77700}
\definecolor{notebg}{HTML}{EAF2FA}

\titleformat{\section}{\sffamily\Large\bfseries\color{accent}}{\thesection}{0.8em}{}
\titleformat{\subsection}{\sffamily\large\bfseries}{\thesubsection}{0.8em}{}
\titlespacing*{\section}{0pt}{2.2ex plus .5ex}{1.2ex}

\setlist{itemsep=0.25em, topsep=0.4em}
\setlength{\parindent}{0pt}
\setlength{\parskip}{0.55em}

\lstset{
  basicstyle=\ttfamily\small,
  backgroundcolor=\color{codebg},
  frame=single, framerule=0pt, rulecolor=\color{codebg},
  xleftmargin=0.6em, xrightmargin=0.6em,
  aboveskip=0.6em, belowskip=0.6em,
  breaklines=true, columns=fullflexible, keepspaces=true,
  showstringspaces=false, upquote=true,
  literate={~}{{\texttildelow}}1
}

\newtcolorbox{note}[1][Note]{colback=notebg, colframe=accent, boxrule=0.5pt,
  arc=2pt, left=6pt, right=6pt, top=4pt, bottom=4pt,
  fonttitle=\sffamily\bfseries\small, title=#1}
\newtcolorbox{caution}[1][Caution]{colback=warnbg, colframe=warnframe, boxrule=0.5pt,
  arc=2pt, left=6pt, right=6pt, top=4pt, bottom=4pt,
  fonttitle=\sffamily\bfseries\small, title=#1}

\newcommand{\file}[1]{\texttt{#1}}
\newcommand{\cmd}[1]{\texttt{#1}}

\pagestyle{fancy}
\fancyhf{}
\fancyhead[L]{\sffamily\small academic-pdf-reading}
\fancyhead[R]{\sffamily\small Installation guide, v1.2.0}
\fancyfoot[C]{\sffamily\small\thepage}
\renewcommand{\headrulewidth}{0.4pt}

\begin{document}

\begin{center}
  {\sffamily\bfseries\LARGE\color{accent} academic-pdf-reading}\\[0.5em]
  {\sffamily\Large Installation and setup guide}\\[0.8em]
  {\small Skill version 1.2.0 \quad$\cdot$\quad 2 October 2026}
\end{center}

\vspace{0.5em}

\section*{What this skill does}

\textbf{academic-pdf-reading} is a Claude skill for reading academic papers
and textbooks without transcription errors. Text extracted from scientific
documents fails silently: a superscript $-1$ becomes a subtraction, an
overbar vanishes, a prime turns into the digit~0. The skill enforces one
rule --- \emph{text layers are for locating and prose; mathematics is read
from page images} --- and supplies the procedures and a measurement script
that make the rule practical. It also maintains a small per-project file
recording verified facts about that project's documents.

\begin{note}[New in v1.2.0: renditions]
The same uploaded document can be mounted in different forms in different
sessions --- for example as the original PDF in one session and as a ZIP of
page images plus text layers in another --- without you changing anything.
Earlier versions reported every such switch as a changed file. From v1.2.0
the inventory records several \emph{renditions} per file, the staleness
check recognises any recorded one, and an unrecorded form of an apparently
unchanged document is reported as a \emph{probable rendition swap} that
Claude confirms with you instead of a full re-triage
(Sections~\ref{sec:daily} and~\ref{sec:swap}).
\end{note}

\begin{note}[Scope of this guide]
This guide covers installing the skill, setting up a project, and keeping
the setup current. Menu names and locations in the Claude app change over
time; where this guide refers to a setting, check
\url{https://support.claude.com} for its current location.
\end{note}

\tableofcontents

\section{The components}

There are three deliverables. Only the first is required.

\begin{center}
\begin{tabularx}{\linewidth}{@{}>{\ttfamily}p{5.4cm}X@{}}
\toprule
\normalfont\textbf{File} & \textbf{Purpose} \\
\midrule
academic-pdf-reading-\allowbreak v1.2.0.zip (or .skill) & The installable skill: a ZIP
of the skill folder. A \file{.skill} file is the same ZIP under another
extension. Install once; available in every chat and project. \\
academic-pdf-reading.md & The same skill as one Markdown file, for use
\emph{without} installing (upload into a project's knowledge). Generated
from the skill; never edit it by hand. \\
pdf\_protocol\_<project>.md & One per project: measured facts about that
project's files, plus a machine-readable inventory. Generated by Claude
with the skill, uploaded by you. \\
\bottomrule
\end{tabularx}
\end{center}

Inside the skill:

\begin{lstlisting}
academic-pdf-reading/
  SKILL.md                         policy, triage, renditions, access
                                   recipes, failure modes, quoting rules
  scripts/triage.py                deterministic measurements
  scripts/build_bundle.py          regenerates the .md bundle
  references/project_protocol.md   how to create/update a project protocol
\end{lstlisting}

\section{Requirements}

\begin{itemize}
  \item \textbf{Skills enabled} for your account. On Team and Enterprise
  plans an administrator may need to allow them.
  \item \textbf{Code execution and file creation enabled}, since the skill
  runs \file{triage.py}. Without it the skill's text still applies, but no
  measurements, crops or staleness checks can run.
  \item Nothing to install yourself. The script needs Python~3.8+ and uses
  poppler-utils and Pillow when present; Claude's sandbox provides them, and
  the script degrades gracefully where a tool is missing
  (\cmd{triage.py tools} reports what exists).
\end{itemize}

\section{Installing or updating the skill}

\subsection{Option A (recommended): upload the skill ZIP}

\textbf{First installation.}
\begin{enumerate}
  \item Download the skill ZIP (or \file{.skill}) file.
  \item Install it. When Claude presents a \file{.skill} file in chat, its
  file card normally offers a \textbf{Save skill} button. Otherwise go to
  \textbf{Customize $\rightarrow$ Skills}, click \textbf{+}, choose
  \textbf{Create skill}, then \textbf{Upload a skill}, and select the ZIP.
  \item Make sure the skill is toggled on.
\end{enumerate}

\textbf{Updating an installed version.} The skill keeps the same name
(\texttt{academic-pdf-reading}) across versions.
\begin{enumerate}
  \item In \textbf{Customize $\rightarrow$ Skills}, open the existing
  skill's three-dot menu and look for an \textbf{Upload and replace} option.
  If it exists, use it with the new ZIP.
  \item If it does not, delete the old skill and upload the new ZIP as for a
  first installation. Nothing is lost: the ZIP contains the complete skill,
  and project protocols live in project knowledge, not in the skill.
  \item Confirm that exactly one copy is listed and enabled. Two copies with
  the same name and description compete for the same requests.
\end{enumerate}

\begin{caution}[If the skill is not listed under Skills]
A skill can also arrive through a plugin rather than a direct upload; it is
then not necessarily editable from the Skills list. Update it the way it was
installed, and ask Claude where the skill is mounted if you are unsure
(\file{/mnt/skills/user/} for uploaded skills; other paths for others).
\end{caution}

\textbf{Verify.} In a new chat, ask:
\begin{lstlisting}
Which version of the academic-pdf-reading skill is installed? Print
metadata.version from SKILL.md and TRIAGE_VERSION from triage.py.
\end{lstlisting}
Both should read \texttt{1.2.0}. To check that the skill works, attach any
paper and ask Claude to triage it with the skill: it should run
\cmd{triage.py profile} and report the file's container (genuine PDF, ZIP
archive, or text extraction), its channels, and any text-layer warnings.

\subsection{Option B: use the Markdown bundle without installing}

Use this where skills are unavailable, or to pin a version in one project.

\begin{enumerate}
  \item Upload \file{academic-pdf-reading.md} to the project's knowledge.
  \item In each session Claude extracts the script from Part~3 of the
  bundle. The bundle's opening snippet does this:
\begin{lstlisting}[language=Python]
import re
t = open('/mnt/project/academic-pdf-reading.md', encoding='utf-8').read()
s = re.search(r'~~~~python\n(.*?)\n~~~~', t, re.S).group(1)
open('/home/claude/triage.py', 'w', encoding='utf-8').write(s)
\end{lstlisting}
  \item Tell Claude to follow the bundle (see the pointer text in
  Section~\ref{sec:pointer}, replacing ``skill'' with
  ``\file{academic-pdf-reading.md} in project knowledge'').
\end{enumerate}

\begin{caution}
If the bundle is uploaded into a project that also has a protocol file, the
staleness check treats the bundle as a source file and does \emph{not} flag
it (from v1.1.0). In v1.0.0 it was flagged as an added file.
\end{caution}

\section{Setting up a project}

Do this once per project whose documents you will ask about repeatedly.
For a one-off paper, skip it: the skill triages files on the fly.

\subsection{Generate the project protocol}

Upload the project's papers first. Then ask:
\begin{lstlisting}
Generate a project-specific PDF protocol for this project, following the
academic-pdf-reading skill's references/project_protocol.md.
\end{lstlisting}
Claude profiles every file, verifies glyph mappings and page offsets on the
page images, and writes \file{pdf\_protocol\_<project>.md} with an embedded
inventory: for each file, the rendition currently mounted (size, SHA-256
hash, container, page count, and a prose sketch used to recognise the same
document in another form later). This is the most token-intensive step, and
it happens once. Other renditions are added as they appear
(Section~\ref{sec:swap}); there is no need to wait for uploads to
``settle'' first.

\subsection{Upload it}

Claude cannot write into project knowledge, so download the file Claude
presents and upload it to the project's knowledge yourself.

\Needspace{12\baselineskip}
\subsection{Point the project instructions at it}\label{sec:pointer}

Add this to the project's instructions (custom instructions). The last
sentence keeps the core rule in force even in a session where the skill
fails to load:

\begin{tcolorbox}[colback=codebg, colframe=codebg, arc=2pt, left=6pt, right=6pt]
\small\raggedright
PDF reading: follow the academic-pdf-reading skill. Project-specific facts
are in \texttt{pdf\_protocol\_<project>.md} in project knowledge; run its
staleness check on first use each session. Invariant even without the
skill: text layers are for locating and prose only; equations, signs,
indices, bars and primes are read from page images, and math from text-only
files is labelled unverified.
\end{tcolorbox}

If the instructions previously contained a full reading protocol, remove it
\emph{after} the next step succeeds; keeping both copies leads to drift.

\subsection{Confirm the read-back}

In a \emph{new} chat in the project, ask:
\begin{lstlisting}
Run the PDF protocol staleness check.
\end{lstlisting}
Expected result: the skill triggers, finds the protocol, and reports
\textbf{up to date}. Only then delete any old protocol text from the
instructions.

\section{Day-to-day use}\label{sec:daily}

The staleness check runs once per session at the first use of the skill,
and again if you say files changed --- not per question. It has three
outcomes:

\begin{center}
\small
\begin{tabularx}{\linewidth}{@{}p{3.6cm}p{1.2cm}X@{}}
\toprule
\textbf{Result} & \textbf{Exit} & \textbf{What Claude does} \\
\midrule
up to date & 0 & Uses the protocol's facts for the rendition that is
mounted (reported as \texttt{mounted} lines when a file has several). \\
\texttt{SWAP?} only (probable rendition swap) & 3 & Lists the files and
asks whether any was edited or replaced. Until you answer, reads them
without the protocol. If you confirm nothing changed, offers to record the
new renditions (Section~\ref{sec:swap}). \\
STALE & 1 & Lists what was added, removed or changed, reads those files
without the protocol, and offers an update. \\
\bottomrule
\end{tabularx}
\end{center}

No outcome blocks an answer.

\begin{itemize}
  \item \textbf{When a file counts as a probable swap.} Its hash matches no
  recorded rendition, but its container differs from every recorded one,
  its page count matches, and its prose sketch (where one is recorded)
  agrees with the recorded document. Anything else is reported as
  \texttt{CHANGED}, including a new hash in an \emph{already recorded}
  container, which may be a revision.
  \item \textbf{Why your answer is needed.} The sketch identifies the
  \emph{document}, not the \emph{version}: a revision that only changes a
  few equations or sentences looks identical to it. Only you know whether a
  file was edited. If you are unsure, say so, and Claude treats the file as
  changed.
  \item \textbf{Source files do not trigger updates.} Adding, editing or
  removing \file{.tex}, \file{.md}, \file{.bib}, \file{.sty} or \file{.cls}
  files is reported as \texttt{source +/-/\textasciitilde} and the result
  stays up to date, so frequent chat-summary uploads cause no update rounds.
  A \file{.txt} file is \emph{not} treated as source: it may be a text
  extraction of a book.
\end{itemize}

\textbf{Keeping token use low.} Ask about specific equations or pages
(``Kugler Eq.~(46)'') rather than whole sections; start a fresh chat per
topic, since images accumulate in context; say when a text-layer reading is
good enough.

\section{Updating the protocol}

\subsection{After a stale result}

When you accept an update, Claude re-triages only the affected files. A
changed file is treated as new: its old renditions and every fact recorded
for them are retired. Claude then regenerates the inventory (keeping the
recorded renditions of all other files), adds a changelog line and hands
you the new file to re-upload.

\subsection{After a confirmed rendition swap}\label{sec:swap}

Once you confirm that the files were not edited, Claude:
\begin{enumerate}
  \item triages the newly mounted rendition and records its
  container-dependent facts (access recipe, glyph mappings, line endings,
  image geometry) in a separate block, leaving the facts of the existing
  rendition untouched;
  \item confirms the document-level facts (page offset, equation-tag range,
  structure) on a few pages of the new rendition;
  \item appends the new rendition to the inventory for exactly the files you
  confirmed, and hands you the protocol to re-upload.
\end{enumerate}
From then on, either rendition passes the check silently, and Claude uses
whichever rendition's facts match the file that is mounted. Mappings are
never carried from one rendition to another: a text layer extracted from a
ZIP and one extracted from the original PDF damage mathematics differently.

\section{Updating the skill itself}

The skill folder is the single source of truth.

\begin{enumerate}
  \item Edit files in the skill folder (or ask Claude to; Claude works on a
  copy, since installed skills are read-only).
  \item In \file{SKILL.md}, bump \texttt{metadata.version} and
  \texttt{metadata.date}; also bump \texttt{TRIAGE\_VERSION} in \file{triage.py} if
  the script changed.
  \item Rebuild the bundle and repackage:
\begin{lstlisting}
python3 academic-pdf-reading/scripts/build_bundle.py academic-pdf-reading.md
python3 -m scripts.package_skill academic-pdf-reading OUTDIR   # skill-creator
\end{lstlisting}
  A plain ZIP of the folder also works:
\begin{lstlisting}
zip -r academic-pdf-reading-<version>.zip academic-pdf-reading
\end{lstlisting}
  \item Reinstall the skill (Section~3.1); re-upload the bundle wherever it
  is used. The bundle's first line carries the version, so a mismatch with
  the installed skill is visible.
\end{enumerate}

\textbf{Protocol compatibility.} v1.2.0 introduces inventory format~2
(several renditions per file). It reads format~1 inventories (v1.0.0 and
v1.1.0), each file becoming a single rendition, so existing protocols keep
working without regeneration. The next protocol update writes format~2.

\section{Troubleshooting}\label{sec:trouble}

\begin{description}[style=nextline, leftmargin=1em, font=\bfseries\sffamily]

\item[The skill does not trigger.]
Name it explicitly (``use the academic-pdf-reading skill''). Check that
it is installed and enabled, and that code execution is enabled.

\item[The check reports \texttt{SWAP?} although I changed nothing.]
Expected after a rendition switch: the platform mounted a form of the file
that the protocol has not recorded yet. Confirm that nothing changed and
accept the offered update; afterwards both forms pass. Switches have been
observed in both directions (PDF $\rightarrow$ ZIP and ZIP $\rightarrow$
PDF), so waiting does not prevent them; recording both renditions does.

\item[The check reports \texttt{CHANGED} although I changed nothing.]
The file's form was not recognisable as the same document: a different
page count, a disagreeing sketch, or a new hash in an already recorded
container (for example, the platform re-rendered a ZIP). Tell Claude
nothing changed; Claude still re-triages the file, because a recorded
rendition's facts cannot be assumed to hold for a new rendering.

\item[A \file{.pdf} is not a PDF.]
Expected. Project files may be ZIP archives (page images plus text layers)
or plain-text extractions with a \file{.pdf} extension. The skill identifies
files by their first bytes, never by extension.

\item[A textbook has no page images.]
Text-only extractions have no image channel. Mathematics from them is
labelled \emph{text-layer reading, unverified} and should be cross-checked
against a paper that does have images, or re-derived.

\item[The bundle's script extraction fails.]
The snippet assumes \file{/mnt/project/academic-pdf-reading.md}; change the
path if the bundle lives elsewhere.

\item[A tool is missing.]
Run \cmd{triage.py tools}. Without a PDF renderer, genuine PDFs fall into
the unverified regime until one is available; Claude will say so.

\end{description}

\appendix
\section{Command reference: \file{triage.py}}

Run from the skill directory as \cmd{python3 scripts/triage.py <command>}.
Exit codes: 0 success or up to date; 1 stale; 2 usage or I/O error;
3 only probable rendition swaps.

\begin{center}
\small
\begin{tabularx}{\linewidth}{@{}>{\ttfamily}p{5.2cm}X@{}}
\toprule
\normalfont\textbf{Command} & \textbf{What it does} \\
\midrule
tools & Lists available helper tools. \\
profile FILE... [-{}-json OUT] & Container, channels, text-layer signals,
page-offset candidates, notes. Measures; never decides. \\
inventory FILE... [-{}-root DIR] [-{}-out OUT] & Inventory for a protocol:
per file, its renditions (size, SHA-256, container, pages, prose
sketch). \\
\quad[-{}-merge OLD] & Keeps OLD's recorded renditions for files whose
hash is already known. A file with a new hash replaces its entry. \\
\quad[-{}-add-rendition NAME...] & With \texttt{-{}-merge}: the named files'
new hashes are appended as further renditions instead. Only for files the
user confirmed unchanged. \\
check INVENTORY [-{}-root DIR] [-{}-json] & Staleness check against a
\file{.json} or a protocol \file{.md}. A file matching any recorded
rendition is unchanged. Exit 1 if stale, 3 if only probable swaps. \\
text FILE [N] [-{}-escape] & One page of a text layer (or a line range of
a text file); \texttt{-{}-escape} shows control bytes. \\
grep PATTERN FILE... [-i] & Regex search reported as file:page:line;
CRLF-safe. \\
extract FILE -{}-out DIR & Unpacks a ZIP; splits a PDF into per-page
\file{N.txt}. \\
crop SRC -{}-page N [-{}-box x0 y0 x1 y1] -{}-out PNG & Page image,
optionally cropped and upscaled, from a ZIP, PDF or image. \\
\bottomrule
\end{tabularx}
\end{center}

\end{document}
